\documentclass[a4paper,12pt]{report}

\usepackage[lmargin=2.00000cm,rmargin=2.0000cm,tmargin=5.5000cm,bmargin=2.500000cm,headheight=3.5cm]{geometry}        %Flexible and complete interface to document dimensions
\usepackage[utf8]{inputenc}
\usepackage[T1]{fontenc}
%\usepackage[latin1]{inputenc}
\usepackage{amsmath}
\usepackage{amsfonts}
\usepackage{amssymb}
\usepackage{lmodern}
\usepackage{float}
\usepackage{graphicx}

\usepackage[english,french]{babel}        %use for the below package `datetime'
\usepackage[babel=true]{csquotes}
\usepackage{datetime}        %Change format of `\today' with commands for current time
\renewcommand{\dateseparator}{-}
\newcommand{\headertoday}{\twodigit\day \dateseparator \twodigit\month  \dateseparator \the\year}

\usepackage{animate}
\usepackage{lastpage}
\usepackage{array}
\usepackage{lastpage}
\usepackage{multirow}
\usepackage{titling}
\usepackage{placeins}
\usepackage{media9}
\usepackage{eurosym}
\usepackage{pdflscape}
\usepackage{color}
\usepackage[table]{xcolor}
\definecolor{mon_bleu}{rgb}{0.137,0.466,0.741}
\definecolor{mon_vert}{rgb}{0.07843,0.4627,0.07843}
\definecolor{mon_rouge}{rgb}{0.62745,0.16078,0.27058}		
%\usepackage{subfigure}
\usepackage{subcaption}

\usepackage[ 
           hidelinks,
           colorlinks=true,
           linkcolor=blue,          % color of internal links (change box color with linkbordercolor)
    	   citecolor=blue,        % color of links to bibliography
    	   filecolor=blue,      % color of file links
    	   urlcolor=blue,           % color of external links
           pdfhighlight =/O]{hyperref}																		% dvipdfm package pour lien hypertextes dans pdf (pdfhighlight--> afficher la main dans le pdf, colorlinks--> colorlinks=true afficher les liens en couleurs)
  
\usepackage{blindtext}
\usepackage[final]{pdfpages}
\usepackage[french]{cleveref}																				% à placer après hyperref attention ne pas utiliser ":" dans les labels des équations avec French babel activé et cref

\usepackage{fancyhdr}
%%%%%%%%%%%%%%%%%%%%%%%%%%%%%%%%%%%%%%%%%%%%%%%%%%%%%%%%%%%%%%%%%%%%%%%%%%
%                         FORMAT PERSONNALISES                           %
% %%%%%%%%%%%%%%%%%%%%%%%%%%%%%%%%%%%%%%%%%%%%%%%%%%%%%%%%%%%%%%%%%%%%%%%%

\parindent=0pt        %leading space for paragraphs
\pagestyle{fancy}
\renewcommand{\arraystretch}{1.5}
\renewcommand{\headrulewidth}{0pt}
\fancyhead[CE,CO,LE,LO,RE,RO]{} %% clear out all headers
\fancyhead[C]{%
\begin{tabular}{|m{3.0cm}|m{10cm}|c@{}|}
\hline
\multirow{2}{*}{\includegraphics[scale=0.045]																			% LOGO LABO
{Logo-Symme.png}}& \centering \multirow{2}{*}{ \Large{\thetitle}}  &  \multirow{2}{*}{Page: \thepage ~/ \pageref{LastPage}}\\
& &  \\
\hline
%Etabli par: \newline \theauthor & \multicolumn{2}{l|}{Diffusion: interne}\\
%\hline
\multicolumn{2}{|l|}{Objet: \thetitleobject }& Date: \headertoday ~~\\ 
\hline
\end{tabular}
}
\renewcommand\footrulewidth{1pt}
\fancyfoot[L]{Document confidentiel} 																		% DOCUMENT CONFIDENTIEL
\fancyfoot[R]{Laboratoire SYMME}		 																	% NOM DU LABO

\newcolumntype{x}[1]{>{\centering\hspace{0pt}}p{#1}}
\setlength{\doublerulesep}{\arrayrulewidth} 


\usepackage{tikz}																							% pour les diagrammes
\usetikzlibrary{calc}
\usetikzlibrary{shapes.geometric,shapes.arrows,decorations.pathmorphing}
\usetikzlibrary{matrix,chains,scopes,positioning,shapes,arrows,fit}
\usetikzlibrary{calc,decorations.pathreplacing}
\usetikzlibrary{calc}
\usetikzlibrary{backgrounds,decorations.markings}

\setcounter{secnumdepth}{5}		%sous sous sous section



%%%%%%%%%%%%%%%%%%%%%%%%%%%%%%%%%%%%%%%%%%%%%%%%%%%%%%%%%%%%%%%%%%%%%%%%%%%%%%%%%%%%
%%-------------> PAGE DE GARDE INFO
%%%%%%%%%%%%%%%%%%%%%%%%%%%%%%%%%%%%%%%%%%%%%%%%%%%%%%%%%%%%%%%%%%%%%%%%%%%%%%%%%%%%

\author{Auteur}
\newcommand{\validator}{J. COLLOMB}
\title{Rapport d'avancement}
\selectlanguage{french}	
\date{\today}
\newcommand{\thetitleobject}{Comment insérer des équations ?}
\setcounter{tocdepth}{6}
\setcounter{secnumdepth}{6}


%%%%%%%%%%%%%%%%%%%%%%%%%%%%%%%%%%%%%%%%%%%%%%%%%%%%%%%%%%%%%%%%%%%%%%%%%%%%%%%%%%%%
%%-------------> DEBUT DU DOCUMENT 
%%%%%%%%%%%%%%%%%%%%%%%%%%%%%%%%%%%%%%%%%%%%%%%%%%%%%%%%%%%%%%%%%%%%%%%%%%%%%%%%%%%%

\begin{document}

\graphicspath{{Figures/}}

%%%%%%%%%%%%%%%%%%%%%%%%%%%%%%%%%%%%%%%%%%%%%%%%%%%%%%%%%%%%%%%%%%%%%%%%%%%%%%%%%%%%
%%-------------> PAGE DE GARDE 
%%%%%%%%%%%%%%%%%%%%%%%%%%%%%%%%%%%%%%%%%%%%%%%%%%%%%%%%%%%%%%%%%%%%%%%%%%%%%%%%%%%%
\begin{titlepage}
    \centering
    \vspace*{3cm}
    {\bfseries\Large
        \thetitle\\
        \thetitleobject\\
        --- \\
        Document confidentiel\\  																			% LIGNE "CONFIDENTIEL"
        \vskip2cm
    }
    \includegraphics[scale=0.65]{logos.png}  																% LOGOS DES PARTENAIRES
    \vfill
    \begin{center}
    \begin{table}[b]
    \begin{tabular}{x{.225\linewidth}|| x{.225\linewidth}|| x{.225\linewidth} || x{.225\linewidth} }
    \hline \hline
    \textbf{Date} & \textbf{Révision} & \textbf{Rédigé par} & \textbf{Validé par} \tabularnewline
    \hline \hline
    \thedate & Rev. A  &  \theauthor &  \validator \tabularnewline
    \hline
    ~ & ~ & ~ & ~ \tabularnewline
    \hline
    ~ & ~ & ~ & ~ \tabularnewline
    \hline
    ~ & ~ & ~ & ~\tabularnewline         
    \hline \hline                                                                   
    \end{tabular}
    \end{table}
    \end{center}
    \vfill
    \vfill
    \begin{tikzpicture}[remember picture,overlay]
    \node (label) at (8cm,20cm){
        \includegraphics[width=2.5cm]{Logo-Symme.png} 													% LOGO LABO
      };
    
	\draw[very thick]
		([yshift=-25pt,xshift=25pt]current page.north west)--
		([yshift=-25pt,xshift=-25pt]current page.north east)--
		([yshift=25pt,xshift=-25pt]current page.south east)--
		([yshift=25pt,xshift=25pt]current page.south west)--cycle;
    \end{tikzpicture}
\end{titlepage}


%\begin{titlepage}
%    \centering
%    
%    \vspace*{3cm}
%    {\bfseries\Large
%        \thetitle\\
%        \thetitleobject\\
%        \vskip2cm
%    }    
%    \vfill
%    \begin{center}
%    \begin{table}[b]
%    \begin{tabular}{x{.225\linewidth}|| x{.225\linewidth}|| x{.225\linewidth} || x{.225\linewidth} }
%    \hline \hline
%    \textbf{Date} & \textbf{R�vision} & \textbf{R�dig� par} & \textbf{Valid� par} \tabularnewline
%    \hline \hline
%    \thedate & Rev. A  &  \theauthor &  \validator \tabularnewline
%    \hline
%    ~ & ~ & ~ & ~ \tabularnewline
%    \hline
%    ~ & ~ & ~ & ~ \tabularnewline
%    \hline
%    ~ & ~ & ~ & ~\tabularnewline         
%    \hline \hline                                                                   
%    \end{tabular}
%    \end{table}
%    \end{center}
%    \vfill
%    \vfill
%    \begin{tikzpicture}[overlay,remember picture]
%    \node (label) at (6.7cm,19.5cm){
%        \includegraphics[width=4cm]{../Figures/logo_CT1.pdf} % also works with logo.pdf
%      };
%    
%    \draw [line width=1pt,rounded corners=7pt]
%        ($ (current page.north west) + (1.5cm,-1.5cm) $)
%        rectangle
%        ($ (current page.south east) + (-1.5cm,1.5cm) $);
%        
%%	\draw[very thick]
%%		([yshift=-25pt,xshift=25pt]current page.north west)--
%%		([yshift=-25pt,xshift=-25pt]current page.north east)--
%%		([yshift=25pt,xshift=-25pt]current page.south east)--
%%		([yshift=25pt,xshift=25pt]current page.south west)--cycle;
%    \end{tikzpicture}
%\end{titlepage}


%%%%%%%%%%%%%%%%%%%%%%%%%%%%%%%%%%%%%%%%%%%%%%%%%%%%%%%%%%%%%%%%%%%%%%%%%%%%%%%%%%%%
%%-------------> SOMMAIRE
%%%%%%%%%%%%%%%%%%%%%%%%%%%%%%%%%%%%%%%%%%%%%%%%%%%%%%%%%%%%%%%%%%%%%%%%%%%%%%%%%%%%

\renewcommand\contentsname{Sommaire}
\setcounter{chapter}{1}
\tableofcontents
%\listoffigures
%\listoftables


%%%%%%%%%%%%%%%%%%%%%%%%%%%%%%%%%%%%%%%%%%%%%%%%%%%%%%%%%%%%%%%%%%%%%%%%%%%%%%%%%%%%
%%-------------> CORPS DOCUMENT
%%%%%%%%%%%%%%%%%%%%%%%%%%%%%%%%%%%%%%%%%%%%%%%%%%%%%%%%%%%%%%%%%%%%%%%%%%%%%%%%%%%%

%-----------------------------------------------------------------------------------
%-----------------------------------------------------------------------------------
\newpage

\section{Équations en ligne}
Sous \LaTeX, il est possible d'intégrer des équations en ligne, c'est à dire dans le texte. Ces équations ne sont pas numérotés. Pour exemple, voici une équation en ligne : $a \times b~=~a.b$. \\

Pour intégrer une équation en ligne, la commande \LaTeX est : \verb|$ équation $|. \\

Cette commande permet d'intégrer :
\begin{itemize}
\item des équations \verb|$ a \times b = a.b $| $ \longrightarrow a \times b = a.b$;
\item des symboles \verb|$ \Delta $, $ \epsilon $, $ \phi $| $ \longrightarrow  \Delta $, $ \epsilon $, $ \phi $;
\item des indices \verb|$ \Delta T_{maximal} $| $ \longrightarrow  \Delta T_{maximal} $;
\item des exposants \verb|$ \Delta T^{maximal} $| $ \longrightarrow  \Delta T^{maximal} $;
\item des fractions \verb|$ \frac{a}{b} $| $ \longrightarrow  \frac{a}{b} $; 
\item des racines carrées \verb|$ \sqrt{abc} $| $ \longrightarrow  \sqrt{abc} $; 
\item des racines $n^{ième}$ \verb|$ \sqrt[n]{abc} $| $ \longrightarrow  \sqrt[n]{abc} $; 
\item \ldots
\end{itemize}


\section{Équations hors ligne}
Il est également possible d'intégrer des équations hors-ligne, c'est-à-dire séparées du texte. Ces équations sont numérotés. Les équations~\ref{eq_densite_flux_conduction} et \ref{eq_equatio_chaleur} sont des équations hors-ligne. \\

Pour intégrer une équation hors-ligne, la commande \LaTeX est : \\
\verb|\begin{equation}| \\
\verb|\overrightarrow{\varphi_{conductif}} = -\lambda(T) \times \overrightarrow{grad(T)}| \\
\verb|\label{eq_densite_flux_conduction}| \\
\verb|\end{equation}|

\begin{equation}
\overrightarrow{\varphi_{conductif}} = -\lambda(T) \times \overrightarrow{grad(T)}
\label{eq_densite_flux_conduction}
\end{equation}

\verb|\begin{equation}| \\
\verb|\rho (T) \times c_{p} (T) \times \frac{\partial T}{\partial t} = | \\
\verb|\frac{\partial }{\partial x} \times (\lambda (T) \times | \\
\verb|\frac{\partial T}{\partial x}) + \frac{\partial }{\partial y} \times | \\
\verb|(\lambda (T) \times \frac{\partial T}{\partial y}) + | \\
\verb|\frac{\partial }{\partial z} \times (\lambda (T) \times | \\
\verb|\frac{\partial T}{\partial z}) + \dot{E}_{generee}| \\
\verb|\label{eq_equatio_chaleur}| \\
\verb|\end{equation}|

\begin{equation}
\rho (T) \times c_{p} (T) \times \frac{\partial T}{\partial t} = \frac{\partial }{\partial x} \times (\lambda (T) \times \frac{\partial T}{\partial x}) + \frac{\partial }{\partial y} \times (\lambda (T) \times \frac{\partial T}{\partial y}) + \frac{\partial }{\partial z} \times (\lambda (T) \times \frac{\partial T}{\partial z}) + \dot{E}_{generee}
\label{eq_equatio_chaleur}
\end{equation}



\section{Sous équations}
\LaTeX permet également la création de sous-équations. Ces équations sont numérotés. Les équations~\ref{eq_equation_singularite_A1} à \ref{eq_equation_singularite_A3}. \\

Pour intégrer des sous-équations, la commande \LaTeX est : \\
\verb|\begin{subequations}| \\
\verb|\begin{align}| \\
\verb|A &= 0,9 \times  \sin (\delta)  & & \text{ pour } {\delta \leq 70^{\circ}}| \\
\verb|\label{eq_equation_singularite_A1}| \\
\verb|\\| \\
\verb|A &= 1 & & \text{ pour } {\delta = 90^{\circ}}| \\
\verb|\label{eq_equation_singularite_A2}| \\
\verb|\\| \\
\verb|A &= 0,7 + 0,35 \times  \frac{\delta}{90} & & \text{ pour } {\delta \geq 90^{\circ}}| \\
\verb|\label{eq_equation_singularite_A3}| \\
\verb|\end{align}| \\
\verb|\end{subequations} | \\

\begin{subequations}
\begin{align}
A &= 0,9 \times  \sin (\delta)  & & \text{ pour } {\delta \leq 70^{\circ}}
\label{eq_equation_singularite_A1}
\\
A &= 1 & & \text{ pour } {\delta = 90^{\circ}}
\label{eq_equation_singularite_A2}
\\
A &= 0,7 + 0,35 \times  \frac{\delta}{90} & & \text{ pour } {\delta \geq 90^{\circ}}
\label{eq_equation_singularite_A3}
\end{align}
\end{subequations} 



\section{Matrices}
Il est possible sous \LaTeX d'intégrer des matrices. Par exemple, équation~\ref{matrice} et \ref{matrice2}.

\begin{equation}
\begin{bmatrix}
a & b & c \\ 
d & e & f \\ 
g & h & i
\end{bmatrix} 
\label{matrice}
\end{equation}

\begin{equation}
\begin{pmatrix}
1 & 2 \\ 
3 & 4
\end{pmatrix}
\label{matrice2}
\end{equation}



%%%%%%%%%%%%%%%%%%%%%%%%%%%%%%%%%%%%%%%%%%%%%%%%%%%%%%%%%%%%%%%%%%%%%%%%%%%%%%%%%%%%
%%-------------> FIN DU DOCUMENT
%%%%%%%%%%%%%%%%%%%%%%%%%%%%%%%%%%%%%%%%%%%%%%%%%%%%%%%%%%%%%%%%%%%%%%%%%%%%%%%%%%%%

\end{document}